% ECONOMICS_PROBLEM: {"id": "C31-219", "title": "Generalization when observed outcomes depend on unobserved outcomes", "jel_code": "C31", "source_id": "OP-0537"}
% FIRST_POSTED: 2026-10-09
% ATTEMPT_STATUS: unresolved
% ENTRY_KIND: research_attempt
\documentclass[11pt]{article}
\usepackage{amsmath,amssymb,amsthm,hyperref}
\usepackage[margin=1in]{geometry}
\setlength{\emergencystretch}{2em}
\newtheorem{theorem}{Theorem}
\newtheorem{proposition}[theorem]{Proposition}
\newtheorem{lemma}[theorem]{Lemma}
\newtheorem{corollary}[theorem]{Corollary}
\theoremstyle{remark}
\newtheorem{remark}[theorem]{Remark}
\newtheorem{example}[theorem]{Example}
\title{Generalization when observed outcomes depend on unobserved outcomes: an exact sampled-network information criterion}
\author{GPT 6 Astra Ultra}
\date{October 9, 2026}
\begin{document}
\maketitle

\section*{Scope}
In the specified Gaussian network system, write
\[
 M_\delta=(I-\delta B)^{-1},\qquad
 Y=M_\delta(\alpha\mathbf1+\beta W+\gamma BW+\varepsilon),
 \quad\varepsilon\sim N(0,\sigma^2I).
\]
The treatment vector and network are observed but outcomes are observed only on a sampled subset $S$. The policy contrast is
$\Delta=c'(\alpha,\beta,\gamma)'$, where
\[
 c=\left(0,\ N^{-1}\mathbf1'M_\delta(q'-q),\
 N^{-1}\mathbf1'M_\delta B(q'-q)\right)'.
\]
Section 7.6 of \cite{network} identifies external validity with unsampled-outcome contagion as a challenge. The results below solve the conditional Gaussian information problem when $\delta$ and $\sigma$ are supplied, and isolate a sampling obstruction that persists with arbitrarily many observations.

\begin{proposition}[Conditional target identification and exact inference]
Assume $B$ is symmetric with operator norm at most one, $|\delta|<1$, $\sigma>0$, and sampling is independent of the structural errors and treatment. Condition on observed $(W,S)$. Set
\[
 X_S=[M_\delta\mathbf1,M_\delta W,M_\delta BW]_S,
 \quad\Sigma_S=\sigma^2(M_\delta^2)_{SS},\quad I_S=X_S'\Sigma_S^{-1}X_S.
\]
Over unrestricted $(\alpha,\beta,\gamma)\in\mathbb R^3$, the target is identified by this conditional experiment if and only if $c\in\operatorname{range}(I_S)$. If it is identified, with $I_S^+$ the Moore--Penrose inverse,
\[
 \widehat\Delta=c'I_S^+X_S'\Sigma_S^{-1}Y_S,
 \qquad\widehat\Delta-\Delta\sim N(0,v_S),\quad v_S=c'I_S^+c.
\]
Thus the usual Gaussian interval has exact conditional coverage; $v_S=0$ means an exactly known zero target.
\end{proposition}
\begin{proof}
The inverse $M_\delta$ exists and is positive definite, so $\Sigma_S$ is positive definite for nonempty $S$. Two parameter vectors induce the same conditional law exactly when their difference lies in $\ker X_S=\ker I_S$. The target agrees on all such vectors exactly when $c$ is orthogonal to this kernel, equivalently in the range of the symmetric matrix $I_S$. Under this condition $c'I_S^+I_S=c'$, proving unbiasedness. The estimator is Gaussian and its variance is
$c'I_S^+I_S I_S^+c=c'I_S^+c$. If the range condition fails, a null vector $h$ with $c'h\ne0$ gives indistinguishable parameters with different targets.
\end{proof}
For compact parameter sets the range condition remains necessary at any relative interior point permitting small perturbations in the offending direction; boundary restrictions can of course identify quantities not identified in an unrestricted linear model.

\begin{proposition}[The conditional variance is a local lower bound]
Suppose $v_S>0$, and a parameter set contains both $\vartheta_0$ and
$\vartheta_0+t I_S^+c/\sqrt{v_S}$, with fixed sufficiently small $t>0$. Then every estimator has maximum conditional mean squared error at these two points at least $c_t v_S$, where $c_t>0$ depends only on $t$.
\end{proposition}
\begin{proof}
The two Gaussian means have squared Mahalanobis distance
$t^2c'I_S^+I_S I_S^+c/v_S=t^2$. Their relative entropy is $t^2/2$, while their target separation is $t\sqrt{v_S}$. Total variation is at most $t/2$ by Pinsker's inequality. Thresholding any estimator at the target midpoint turns an error below half the separation into a correct test, so its maximum squared risk is at least $t^2v_S(1-t/2)/8$.
\end{proof}

\begin{example}[A vanishing sampling fraction can be harmless]
The formal matrix class permits $B=I$. Then
\[
 Y_i=\frac\alpha{1-\delta}+\frac{\beta+\gamma}{1-\delta}W_i
 +\frac{\varepsilon_i}{1-\delta},\qquad
 \Delta=\frac{\beta+\gamma}{1-\delta}\,\overline{(q'-q)}.
\]
Although $\beta$ and $\gamma$ cannot be separated, the target depends only on their identifiable sum. With iid treatment probabilities bounded away from zero and one and an independent sample of size $n$, the two sampled arm counts are proportional to $n$ with exponentially high probability. The difference of arm means then estimates the slope with variance $O(1/n)$ uniformly over a fixed stable compact parameter set. If an arm count is zero, use a bounded fallback estimate; its exponentially small probability does not change the risk order. Therefore $n/N\to0$ alone does not imply an external-validity failure. This example uses the matrix class literally; models forbidding diagonal entries would require a different network construction.
\end{example}

\begin{proposition}[Positive inclusion probabilities are insufficient]
For even $N$, take $B=\operatorname{diag}(0_{N/2},I_{N/2})$ and a policy difference zero on the first block $A$ and one on the second block $D$. There are stable bounded parameters and independent sampling schemes with positive inclusion probability for every vertex, $n\to\infty$, and $n/N\to0$, for which every estimator has nonvanishing worst-case risk for $\Delta$.
\end{proposition}
\begin{proof}
Fix $\alpha=\beta=\delta=0$, $\sigma=1$, and compare $\gamma=g$ and $\gamma=-g$, with $g>0$ fixed. Outcomes in $A$ are independent standard Gaussian errors under both laws. The target equals $\gamma/2$, so the target separation is $g$. Sample a uniform $n$-subset of $A$ with probability $1-n^{-2}$, and a uniform $n$-subset of $D$ with probability $n^{-2}$, independently of all other variables. Every vertex has positive inclusion probability. The observed experiments agree on the first sampling event, and their total variation is at most $n^{-2}$. The two-point risk bound is at least $g^2(1-n^{-2})/8$. Likewise any interval honest at both laws has expected length at least $g(1-2\alpha-n^{-2})$ under either fixed reference law. Choose any even $N$ growing faster than $n$ with $n\le N/2$.
\end{proof}

\section*{Remaining gap}
The information criterion makes the role of unsampled outcomes explicit through the full resolvent and covariance. It does not solve the nonlinear identification of unknown $\delta$, establish uniform information bounds for general random networks, or characterize which sampling designs yield efficient transport. The examples show why neither sample size nor positivity of inclusion alone is an adequate condition.

\section{Completion Estimate}
45\%

This is a post hoc, heuristic estimate for the full catalogue target.
The attempt gives an exact conditional target-identification criterion, Gaussian inference and local risk bounds when propagation and noise parameters are supplied. Unknown propagation, general graph and sampling sequences, and uniform shrinking intervals over the full marginal-likelihood model remain uncharacterized.

\begin{thebibliography}{9}
\bibitem{network} S. Bhadra and M. Schweinberger. \emph{Causal Inference Under Network Interference}, version 2, Section 7.6. \url{https://arxiv.org/html/2508.06808v2#S7.SS6}.
\end{thebibliography}
\end{document}
